\begin{abstract}
类型化决策模型返回所声明选项上的概率而非文本，从而使逐条筛查所有消息的成本变得低廉。
本文围绕 PoL 治理，即依据爱2证明（\pol{}）规范进行的治理，探讨此类模型能否充当其保障机制。\pol{} 是一份开放的治理文本，规定了拦截恨语的“安全阀”、仅由 AI 进行的福利争议裁决以及逐人评估。
我们从其条款中推导出七项可检验的要求，并依据首个此类模型 \jev{} 及其开源对应模型在发布后最初 17 天内的 \nPapers{} 项研究和 \nGrey{} 项灰色文献对这些要求进行评估；其中与之相关的 \nCoded{} 项由作者指导的 AI 智能体全文阅读、质量评价并编码（盲法第二编码：$\kappa = \nKappa{}$）。
证据支持在附加条件下将类型化模型用作检测器：它们对风险内容的排序效果良好，成本仅为 LLM 评估器的一小部分；在一项基于模拟数据的预注册研究中，仅当一条独立构建的规则同时认可时才接受其低风险判定，这一做法未造成准确率损失，但伪装攻击仍得以通过。
证据不支持将其用作裁决者：阈值无法迁移，判定可能取决于选项名称，并可能被插入的观点所操纵，“不知道”或“以上皆非”选项的选择不一致，且同类评估器几乎重复了某一类型化模型的全部高置信错误。
类型化模型总体上并不劣于 LLM 评估器，但多数失效缺乏对照比较；截至 10 月 8 日的更新未推翻任何一项回答。
随后，本文对照欧盟法律，指出该文本中的五项张力：情绪识别、二值判定、封闭模型下的透明性、仅由 AI 进行的裁决，以及逐人评分。
我们对 \pol{} 的结论（确定性为低或极低）是：依现有证据，\jev{} 不适合承担任何作出决定的功能，只能如本文所提议的那样，作为有记录、三值、经本地校准的检测器服务于安全阀，且具有实质后果的决定最终由人作出；即便如此，这一结论也是从英文基准外推而来的，因为尚无研究在中文中测量 \pol{} 的构念。
\end{abstract}
